\subsection{Framework and Notation}
\label{sec:frameworknotation}

Let $\epsilon_i$ be the cause indicator such that $\epsilon_i = j$ means individual $i$ experienced failure from cause $j$, where $j = 1,\dots,J$. Our focus is on a priority cause of interest. In cases with more than two causes, we collapse the causes into two such that $Pr(\epsilon = 1) + Pr(\epsilon = 2) = 1$: the PC (cause 1) and all others (cause 2).
%Because our proposed method assumes a priority cause of interest, it is compatible with multiple risks. If there are more than two events are possible, then we collapse the causes into the priority cause and all other causes. Since there are always only two mutually exclusive events, henceforth referred to as cause 1 (priority cause) and cause 2 (all other causes), $Pr(\epsilon = 1) + Pr(\epsilon = 2) = 1$. 
Let $T_{1i}$ and $T_{2i}$ denote the failure times for causes 1 and 2, respectively, for patient $i$, and $T_i = \min(T_{1i}, T_{2i})$ be the overall survival time, representing time to the first failure. Right-censoring due to study end ($\tau$) or loss to follow-up is denoted $C_i$, with censoring time distribution $G_i(t) = \Pr(C_i \geq t | Z_i)$. The observed time is $X_i = \min(T_i, C_i)$. The OS event indicator is $\delta_{i} = I(T_i \le C_i)$. We define $\delta_{1i} = I(T_{1i} \le C_i\wedge T_{2i})$ and $\delta_{2i} = I(T_{2i} \le C_i\wedge T_{1i})$. 

We assume a finite treatment space $\mathcal{A}$. Let $Z_i = (Z_{i1},\dots,Z_{iP})^T$ represent the $P$ covariates. For each $z \in Z$, let $\psi(z) \subseteq A$ to be the set of allowable treatments for a patient with $Z = z$. We observe data composed of $n$ i.i.d replicates $\mathcal{D}_n = \{ Z_i, X_i, \delta_{i}, \delta_{1i} \}_{i=1}^n$. When unambiguous, we drop the index $i$ for simplicity. 

Let $d$ represent an ITR, which maps patient covariate information ($z_i \in Z$) to the treatment assignment ($a_i \in \mathcal{A}$), $d(z_i): z_i \rightarrow a_i$ for each patient $i$, which satisfies $d(z) \in \psi(z)$ for all $z \in Z$.
Let $\phi(\cdot)$ be an optimization criteria for OS and CI. In this manuscript, we focus on the truncated (at $\tau$) area under the curves, $\phi(S) =\int_0^\tau S(t) \, dt$ and $\phi(F_j) = \int_0^\tau F_j(t) \, dt$, the restricted mean survival time at time $\tau$ and integrated CIF truncated at $\tau$ respectively. The latter quantity measures the accumulation of the probability of having the primary event of interest during the follow-up period.

For any treatment $a \in \mathcal{A}$ and covariates vector $Z$, let $S^{(a)}(\cdot \mid Z)$ and $F_j^{(a)}(\cdot \mid Z)$ be the potential survival and cumulative incidence curves for patients with covariates $Z$ and treatment $a$, respectively. Let $\valueonePOT = \phi(S^{(a)}(\cdot \mid Z))$ and $\valuetwoPOT = \phi(F_j^{(a)}(\cdot \mid Z))$ denote the potential outcomes of interest. 
% VALIUE FUNCTION %
%$\mathbb{V}(d(Z))$%
Under regime $d$, define the potential outcomes as $\phi_1^{(d)}(Z) = \sum_{a \in \mathcal{A}} \phi_1^{(a)}(Z) \mathbf{1}\{d(Z) = a\}$ and $\phi_2^{(d)}(Z) = \sum_{a \in \mathcal{A}} \phi_2^{(a)}(Z) \mathbf{1}\{d(Z) = a\}$, where $\mathbf{1}\{B\}$ is the indicator of $B$.
We assume that a prolonged time to failure is optimal, so we optimize a criterion ($\phi_1$) to maximize OS. Additionally, we want to minimize the CIF of the PC via $\phi_2$ due to its inherent importance over all other causes. 
{We aim to estimate the optimal ITR $d_{\text{opt}}$ that optimizes a joint criterion $\boldsymbol{\phi}^{(d)}(Z) = \begin{pmatrix}
\phi_1^{(d)}(Z) \\ 
\phi_2^{(d)}(Z)
\end{pmatrix} = 
\sum_{a \in \mathcal{A}} \begin{pmatrix}
\phi_1^{(a)}(Z) \\ 
\phi_2^{(a)}(Z)
\end{pmatrix}
\cdot \mathbf{1}\{d(Z) = a\}$.
} 
For any treatment regime $d(Z)$, we define the joint value function of the ITR estimator as
\begin{equation*}
    \mathbb{V}(d) = 
    \sum_{a \in \mathcal{A}} \mathbb{E}\left[ 
    \boldsymbol{\phi}^{(a)}(Z)
    %\begin{pmatrix}\valueone \\ \valuetwo\end{pmatrix}
\cdot \mathbf{1}\{d(Z) = a\}
\right].
\end{equation*}
Define $\actiononemax = \argmax_{a \in \mathcal{A}} \valueonePOT$ with $\valueonemax = \phi^{(\actiononemax)}_1(Z)%V_1(a_1^*,Z)
$. If OS across treatments are sufficiently close (i.e., within a tolerance level $\tol$), then for a treatment $a$, 
\begin{equation}
    \label{eqn:TolCondition}
   1 - \frac{\valueonePOT}{\valueonemax} \le \tol.
\end{equation}
This condition assesses if the treatment $a$ performs nearly as well as the best available treatment within the acceptable margin. Equation \cref{eqn:TolCondition} implies $\valueonePOT \ge (1-\tol)\valueonemax$. Thus, the optimal ITR is defined as
$
     \dopt(Z) = 
     \begin{cases}
     \argmin_{a \in \mathcal{A}} \valuetwoPOT & \text{if } \valueonePOT \ge (1-\tol)\valueonemax,\\
     \argmax_{a \in \mathcal{A}} \valueonePOT & \text{otherwise.}
     \end{cases}
$

\subsection{Proposed Estimators}
\label{sec:proposedmethod}

To identify the potential outcomes, we assume Assumptions \ref{assump:Causal1}-\ref{assump:Causal3} \citep{kidwell2016, kosorok_precision_2019}. Then, using identified data, the value function can be written as \\
$\mathbb{V}(d) = 
    \sum_{a \in \mathcal{A}} \mathbb{E}\left[ 
    \boldsymbol{\phi}(a,Z)
\cdot \mathbf{1}\{d(Z) = a\}
\right]$, and we want to estimate the conditional probabilities $\phi_1(a,Z) = \int_0^\tau S(t|Z,A=a) \, dt$ and $\phi_2(a,Z) = \int_0^\tau F_j(t|Z,A=a)\,dt$.

\begin{assumption}[Stable Unit Treatment Value Assumption (SUTVA)]
\label{assump:Causal1}
Each individual's counterfactual outcomes, such as cause-specific failure time or overall failure time, are not affected by the treatment assignments or histories of other individuals in the study. Additionally, the process of censoring is assumed to be independent of other individuals' treatments and histories.
\end{assumption}

\begin{assumption}[Strong ignorability] % look up the name
\label{assump:Causal2}
The treatment assignment is conditionally independent of future potential outcomes (such as $\phi_1^{(a)}$ or $\phi_2^{(a)}$), given the individual's covariates $Z$. This can be expressed as $\phi^{(a)} \perp A \mid Z$, for all $a \in A$. 
\end{assumption}

\begin{assumption}[Positivity]
\label{assump:Causal3}
Given covariate information, the probability of having each treatment $a \in \mathcal{A}$ is greater than a constant $L > 0$, or $\Pr(A = a | Z = z) > L$ for all $a \in \mathcal{A}$, $z$ such that $\Pr(Z = z) > 0$.
\end{assumption}

\subsubsection{Random Forests}
To estimate individualized curves for OS and PC cumulative incidence, we develop a random survival forest (RSF) and random cumulative incidence forest (RCIF), respectively. We adapt the generalized random survival forest in \cite{cho_multi-stage_2023} to the single stage treatment setting. We provide details on the construction of the proposed RSF in the Supplementary Material.

%BELOW IS ABOUT CIF
Following the construction of the RSF, we develop RCIFs to obtain individual CIF curves for optimization (see Algorithm \ref{alg:RCIF}). The Aalen-Johanson (AJ) estimator $\widehat{F}_j$ is used to estimate the cumulative incidence function, as well as for node splitting.
%% below is about aalen-johanson estimator which is used for gray's test and for \phi_2
%For our Random Cumulative Incidence Forest and the CIF criteria $\phi_2$
Letting $k$ represent the node, the AJ estimator for the CIF for cause $j$ failure (i.e., $\epsilon = j$) in group $k$ is given by:
\begin{equation}
\label{AJ_estimator}
   \widehat{F}_{jk}(t) = \intzerotot\widehat{S}_k(u-)Y_k^{-1}(u)dN_{jk}(u) 
\end{equation} where $\widehat{S}_k(t-)$ is the left-hand limit of the Kaplan-Meier estimate $\widehat{S}_k(t)$ for OS (see Equation (S1) in Supplementary Material), $N_{jk}(t)$ is the number of failures of cause $j$ by $t$, $Y_k(t)$ is the number of subjects still at risk just prior to $t$ in group $k$ \citep{gray1988}. In building the trees for the RCIF, $k$ denotes the two groups that are being compared; however, for prediction, $k = 1, \dots, K$ indexes the terminal node, where $K$ is the number of terminal nodes. Whenever $F_{jk}(t)$ refers to the CIF within terminal nodes, we omit the index $k$ such that $F_{jk}(t) \equiv F_j$, since there is no risk of ambiguity. 

% Gray's test for RCIF
In order to partition a node into two potential daughter nodes ($\Node_1, \Node_2$), we use Gray's two-sample test \citep{gray1988} to compare the priority cause CIF between groups: %and calculate the AJ estimator within each node ($\widehat{F}_{jk}$)
\begin{equation}
\label{grays_test}
\text{Gray}(\Node_1, \Node_2) = \int_0^\tau K(t)\{[1-\widehat{F}_{11}(t-)]^{-1}d\widehat{F}_{11}(t)-[1-\widehat{F}_{12}(t-)]^{-1}d\widehat{F}_{12}(t)\},
\end{equation} where $\widehat{F}_{1k}$ is an estimate of $F_{1k}$ based on the AJ estimator for CIF of failure cause $1$ in group $k$ and $K(t)$ is some weight function. In this paper, we define $K(t) = 1$. The splitting rule chooses \( (\Node_1, \Node_2) = \arg \max_{\Node_1,\Node_2} \text{Gray}(\Node_1,\Node_2) \) as the best partition. The recursive partitioning continues with a minimum of $n_{\text{minevent}}$ events (from any cause) in each node while each node's size exceeds the threshold \(n_{\min}\). The PC cumulative incidence estimate of the terminal node \( \Node \) of the \( w \)-th cumulative incidence tree is given by the AJ estimator \( \widehat{F}_{j,\{w\}}(t \mid \Node) \), where \( w = 1, 2, \ldots, n_{\text{tree}} \). Letting $h = (a,Z)$, the RCIF estimate for priority cause $j$ is given by
\[
\widehat{F}_j(t \mid h) = \frac{1}{n_{\text{tree}}} \sum_{w=1}^{n_{\text{tree}}} \sum_{l=1}^{\#\{\Node_{\{w\},l} : l\}} 1(h \in \Node_{\{w\},l}) \widehat{F}_j(t \mid \Node_{\{w\},l}),
\]
where \( \#\{\Node_{\{w\},l} : l\} \) is the number of terminal nodes of the \( w \)-th tree.

%{
%\spacingset{1}
%\begin{spacing}
%%\newcommand{\algorithmOverview}{
\begin{algorithm}[H]
  \caption{The proposed ITR estimator accounting for competing risks.}
  \label{alg:Overview}
  \begin{algorithmic}
  \STATE \textbf{Input:} $Z_i, A_i, \tol$
  \STATE \textbf{Output:} $\dopt$
  \STATE %\ref{alg:RSF}
  \STATE For each individual $i$, given covariates $Z_i$ and small $0<\tol<1$, consider feasible treatments $a_i \in \psi(z_i) 
  \subset \mathcal{A}$ (we drop $i$ for simplicity):
  \begin{enumerate}
    \item Obtain $\survfunchat(\cdot|Z,A)$ via generalized random survival forest \citep{cho_multi-stage_2023}
    \item Obtain $\hat{F}_j(\cdot|Z,A)$ via random cumulative incidence forest (Algorithm \ref{alg:RCIF})
    \item Let optimal treatments for Phase 1 (overall survival) be $\valueonemaxTRUE = \argmax_{a \in \psi(z)}\valueone$
    \item Obtain overall survival treatment set $A^s = \{a \in \psi(Z): \valueone \ge (1- \tol)\valueonemaxTRUE\}$
        \begin{itemize}
            \item If $A^s$ is singleton, $A^s$ is the optimal treatment: $\dopt(Z) = A^s$.
            \item If $\#A^s > 1$, optimal treatment $\dopt(Z) = \argmin_{a \in \textit{A}^s} \valuetwo$.
        \end{itemize} 
    \end{enumerate}
  \STATE In other words, if $a_i \neq A^s_i$ for individual $i$, then individual $i$ stopped at either 
  \begin{itemize*}%[label={Phase \arabic*:}]
      \item \hspace{1em} Phase 1: $\valueone < (1-\tol)\valueonemaxTRUE$ or
      \item \hspace{1em} Phase 2: $\valueone \ge (1-\tol) \valueonemaxTRUE$ but $V_2(a = \dopt, Z) < V_2(a = a_i, Z)$
  \end{itemize*}\\
  and $A_i^s$ may have $\#A_i^s>1$.
\end{algorithmic}
\end{algorithm}
}

\newcommand{\algorithmRSF}{
\begin{algorithm}[H]
  \caption{RSF}
  \label{alg:RSF}
  \begin{algorithmic}
  \STATE \textbf{Input:} $Z_i, A_i$
  \STATE \textbf{Output:} $A_i^*$
  \STATE %\ref{alg:RSF}
  \STATE For each individual $i$, given covariates $Z_i$ and small $0<\tol<1$, consider feasible treatments $a_i \in \psi(z_i) \subset \mathcal{A}$:
  \STATE
  \STATE hi
  \STATE
\end{algorithmic}
\end{algorithm}
}

\newcommand{\algorithmgRCIF}{%

\begin{algorithm}[H]
  \caption{the random cumulative incidence forest estimator}
  \label{alg:RCIF}
  \begin{algorithmic}
  \STATE Result: random cause-j cumulative incidence forest $\hat{F}_j(\cdot|Z,A)$
  \STATE Parameters: the number of trees $n_{\text{tree}}$, the number of minimum terminal node size $n_{\text{min}}$
  \STATE \textbf{Input:} $\{(X_i, \delta_{ji}, Z_i, A_i)\}_{i=1}^n$
  \STATE \textbf{Output:} $\hat{F}_j(\cdot|Z,A=a)$
  \STATE 
  \STATE for \underline{treatment arm} $a \in \mathcal{A}$ \textbf{do} \\
  \STATE \quad for \underline{tree $w = 1, 2, \ldots, n_{\text{tree}}$} \textbf{do} \\
  \STATE \quad \quad Partition the feature space with the input data $\{(I(t\leq X_i), \delta_{ji}, Z_i)\}_{i:a_i=a}$ via Gray's test \cref{grays_test} or MD \citep{cho_multi-stage_2023} until node size $< 2n_{\text{min}}$ \\
  \STATE \quad \quad Obtain the terminal node cumulative incidence probability $\hat{F}^a_{j,\{w\}}(t | Z)$ via Aalen-Johanson estimator \cref{AJ_estimator} \\
  \STATE \quad Obtain the random cumulative incidence forest estimator \\
  \STATE \quad $\hat{F}_j(\cdot | Z, A = a) = \frac{1}{n_{\text{tree}}} \sum_{w=1}^{n_\text{tree}} \hat{F}_{j,\{w\}}^{a}(t | Z)$ \\

\end{algorithmic}
\end{algorithm}
}
%%\algorithmgRCIF
\begin{algorithm}[H]
  \caption{The random cumulative incidence forest estimator}
  \label{alg:RCIF}
  \begin{flushleft}
    \begin{algorithmic}
  \State \textbf{Input:} $\{(X_i, \delta_i, \delta_{ji}, Z_i, A_i)\}_{i=1}^n$
  \State \textbf{Output:} $\widehat{F}_j(\cdot|Z,A=a)$
  \State Result: random cause-j cumulative incidence forest $\widehat{F}_j(\cdot|Z,A)$
  \State Parameters: number of trees $n_{\text{tree}}$, minimum terminal node size $n_{\text{min}}$, minimum number of events $n_{\text{minevent}}$
%  \State 
  \State For \underline{treatment arm} $a \in \mathcal{A}$ \textbf{do} 
  \begin{itemize}[label={}]
  \item For \underline{tree $w = 1, 2, \ldots, n_{\text{tree}}$} \textbf{do}
  \begin{itemize}[label={}]
  \item Partition the feature space with the input data $\{(X_i, \delta_i, \delta_{ji}, Z_i)\}_{i:a_i=a}$ via Gray's test \cref{grays_test} until node size $< 2n_{\text{min}}$ and there are at least $n_{\text{minevent}}$ events from any cause in every node
  \item Obtain the terminal node cumulative incidence probability $\widehat{F}^a_{j,\{w\}}(t | Z)$ via Aalen-Johanson estimator \cref{AJ_estimator}
  \end{itemize}
  
  \item Obtain the random cumulative incidence forest estimator $\widehat{F}_j(\cdot | Z, A = a) = \frac{1}{n_{\text{tree}}} \sum_{w=1}^{n_\text{tree}} \widehat{F}_{j,\{w\}}^{a}(t | Z)$
  \end{itemize}
\end{algorithmic}
\end{flushleft}
\end{algorithm}
%\end{spacing}
%}

\subsubsection{Estimation Procedure}
\label{sec:methodsapproach}
We propose a two-phase estimation procedure that prioritizes maximizing the criterion of OS in Phase 1 but also minimizes the criterion of the CIF of a priority cause in Phase 2 over all treatments for each individual. Whether an individual's treatment is obtained from Phase 1 or Phase 2 depends on a pre-specified tolerance level $0<\tol<1$.
%$\tol$ is defined in Section \ref{sec:methodsapproach}. 

Let $\survfunchat(\cdot;a,Z)$ and $\widehat{F}_j(\cdot;a,Z)$ denote the estimated proposed RSF and RCIF forest estimators, respectively. 
% q-functions
Letting the subscript denote the method phase, we define the true and estimated functions of $\valueone$ and $\valuetwo$ as $\criterionvector^0(a,Z) = \valueTRUE$ and $\widehat{\criterionvector}(a,Z) = \begin{pmatrix} \valueoneHAT \\ \valuetwoHAT \end{pmatrix}$, respectively.
%Remember that $\phi_1$ is overall survival and $\phi_2$ is CIF for the priority cause first mentioned in notation (Section 2.1\label{sec:notation}) then further described in simulations (Section 4.1\label{sec:sims}).
Let the true optimal treatment for Phase 1 be $\actiononemaxTRUE = \argmax_{a \in \mathcal{A}} \valueoneTRUE$ with $\valueonemaxTRUE = V_1^0(\actiononemaxTRUE,Z)$, and the estimated value as $\actiononemaxHAT = \argmax_{a \in \mathcal{A}} \valueoneHAT$ with
$\valueonemaxHAT = \widehat{V}_1(\actiononemaxHAT,Z)$. We can then use the following two-phase procedure to estimate the optimal ITR. For each patient:
\begin{enumerate}
    \item[(i)] estimate the conditional overall survival $\survfunc(t|Z,A)$ for each treatment
    \item[(ii)] estimate the conditional priority j-cause CIF $F_{j}(t|Z,A)$ for each treatment
    \item[(iii)] Phase 1: optimize (i) over feasible finite treatment space for all \(a \in \psi(z) \subset \mathcal{A}\) to obtain optimal treatment for overall survival
    \item[(iv)] Find $A^s$, the set of treatments chosen using $\tol$ criteria to compare to the optimal overall survival treatment
    \begin{enumerate}
        \item if the cardinality of $A^s$ is $1$ then we have ITR estimate $\widehat{d}(z) = A^s$
        \item if the cardinality of $A^s$ is greater than $1$ then optimize (ii) (Phase 2) over treatment set $A^s$ to obtain ITR estimate $\widehat{d}(z)$.
    \end{enumerate}
\end{enumerate}

\begin{algorithm}[H]
  \caption{The proposed ITR estimator accounting for competing risks.}
  \label{alg:Overview}
  \begin{flushleft}
  \begin{algorithmic}
  \State \textbf{Input:} $Z_i, A_i, \tol$
  \State \textbf{Output:} $\doptHAT$
%  \State
  \State For each individual $i$, given covariates $Z_i$ and small $0<\tol<1$, consider feasible treatments $a_i \in \psi(z_i) 
  \subset \mathcal{A}$ (we drop $i$ for simplicity):
  \begin{enumerate}
    \item Obtain $\survfunchat(\cdot|Z,A)$ via random survival forest \citep{cho_multi-stage_2023}
    \item Obtain $\widehat{F}_j(\cdot|Z,A)$ via random cumulative incidence forest (Algorithm \ref{alg:RCIF})
    \item Let the estimated optimal treatments for Phase 1 (overall survival) be $\hat{a}_1^* = \argmax_{a \in \psi(z)}\valueoneHAT$ and $\valueonemaxHAT = \widehat{V}_1(\hat{a}_1^*,Z)$
    \item Obtain overall survival treatment set $A^s = \{a \in \psi(Z): \valueoneHAT \ge (1- \tol)\valueonemaxHAT \}$
        \begin{itemize}
            \item If $A^s$ is singleton, $A^s$ is the estimated optimal treatment: $\doptHAT(Z) = A^s$.
            \item If $\#A^s > 1$, optimal treatment for Phase 2 is $\doptHAT(Z) = \argmin_{a \in \textit{A}^s} \valuetwoHAT$.
        \end{itemize} 
    \end{enumerate}
  \State In other words, if $a_i \neq A^s_i$ for individual $i$, then individual $i$ stopped at either 
    \begin{itemize}[label={}]
      \item \hspace{1em} Phase 1: $\valueoneHAT < (1-\tol)\valueonemaxHAT$ or
      \item \hspace{1em} Phase 2: $\valueoneHAT \ge (1-\tol) \valueonemaxHAT$ but $\valuetwoHATnoarg(a = \doptHAT, Z) < \valuetwoHATnoarg(a = a_i, Z)$
  \end{itemize}
and $A_i^s$ may have $\#A_i^s>1$.
\end{algorithmic}
\end{flushleft}
\end{algorithm}

Algorithm \ref{alg:Overview} provides a high-level outline of our proposed method. Briefly, if the areas under the truncated survival curve between the treatments are within certain $\tol$, then we want to reduce PC cumulative incidence. We consider the competing risk of a priority cause if the OS experience is not too different. However, if the OS experience is excessively different, then we proceed as if the competing risk is not relevant. From Equation \cref{eqn:TolCondition}, $1 - \frac{\valueoneHAT}{\valueonemaxHAT} \le \tol \implies \valueoneHAT \ge (1-\tol)\valueonemaxHAT$. Thus, the optimal ITR when $\#A^s > 1$ is defined as $\doptHAT(Z) = \argmin_{a \in \mathcal{A}:\valueoneHAT \ge (1-\tol)\valueonemaxHAT} \valuetwoHAT$.